\chapter*{Agradecimientos}
~\label{acknoledgements}
\addcontentsline{toc}{chapter}{Agradecimientos}

En primer lugar, quiero agradecer a todas las personas que construyen la
universidad pública, la UNR y el DCC. A todos los profesores de los que tanto
aprendí, haciendo especial mención de Dante, Ceci, Ana, Maite y Mauro, que me
han ayudado a crecer en más de una manera, y a todos mis compañeros que se
volvieron amigos querídisimos: matambre de boga, rio y, con melancolía, pancito
italiano, sepan que fueron la luz en el faro.

En segundo lugar, quiero agradecer a las personas que hicieron que esta tesina
fuera posible. A Pedro por guiarme, enseñarme y tomarse el tiempo de tener
muchas reuniones virtuales y probablemente uno de los hilos más largos de mail,
y a Ceci, por bancarme en aún otra etapa más, con la paciencia y dulzura que la
caracterizan.

En tercer lugar, quiero agradecer a mi familia. A mi mamá, por ser mi sostén
principal, bancarse mis primeros llantos facultativos y siempre inculcarme que
la educación es lo mejor que tenemos al alcance. A mis hermanos: Denise, Eric,
Maggie y Paddy, por enseñarme cada uno a su manera, todo lo que a ellos les
hubiera gustado saber. A Cami, por apoyarme y abrazarme en el proceso. Y a mi
familia que no es de sangre: Frank, Coni, Euge, Mati, Ailo, Vicky y Rampe. No
creo que pueda (ni quiero) dar algún paso importante en la vida sin mencionar a
estas como personas claves en mi desarrollo.

Y en último lugar, a las dos personas que fueron mi todo durante esta
experiencia. Mis compas de \textit{pair-programming} con fracción del teclado
para cada uno, mi audiencia frente a quejas y mi abrazo más esperado en los
logros, los que toman para que yo me empede, las medias neuronas que me
faltaban: Joaquín, quien personifica el DCC, el que siempre va a estar para
darte unos masajitos en la espalda si estás estresado y un abrazo si estás
angustiado, e Inés, la persona con la que viví todo, con la que tenemos la
misma educación, y por lo tanto, la misma visión, la que me acompaña desde los
4 años (y espero que los siga haciendo por siempre). Gracias, a ustedes va
dedicada esta tesina.